\section{Davis--Kahan result inventory}
\label{app:dk-inventory}

\Cref{tab:dk-full-inventory} records the maintained result-level denominator.
Every row is currently verified in the ordinary Lean build and has accepted
semantic/source-correspondence review. The disposition column distinguishes a
proof of the source claim from the checked refutation of Proposition~4.4.

\begin{longtable}{>{\raggedright\arraybackslash}p{0.16\linewidth}>{\raggedright\arraybackslash}p{0.55\linewidth}>{\raggedright\arraybackslash}p{0.18\linewidth}}
\caption{The \DKResultCount{} Davis--Kahan result-level targets.}
\label{tab:dk-full-inventory}\\
\toprule
ID & Result & Disposition \\
\midrule
\endfirsthead
\toprule
ID & Result & Disposition \\
\midrule
\endhead
\\path{S2-sin-theta} & Single-angle sine theorem & proved at source scope \\
\\path{S2-tan-theta} & Single-angle tangent theorem & proved at source scope \\
\\path{S2-sin-two-theta} & Double-angle sine theorem & proved at source scope \\
\\path{S2-tan-two-theta} & Double-angle tangent theorem & proved at source scope \\
\\path{DK-3.1-prop} & Acute direct rotation existence and uniqueness & proved at source scope \\
\\path{DK-3.2-prop} & Nonacute existence criterion & proved at source scope \\
\\path{DK-3.3-prop} & Principal square-root characterization & proved at source scope \\
\\path{DK-3.4-prop} & Square as a direct rotation & proved at source scope \\
\\path{DK-3.1-thm} & Classification of pairs of subspaces & proved at source scope \\
\\path{DK-3.1-cor} & Compact classification by angle eigenvalues & proved at source scope \\
\\path{DK-3.5-prop} & Angle commutation and eigenspace geometry & proved at source scope \\
\\path{DK-3.2-cor} & Reversal symmetry & proved at source scope \\
\\path{DK-4.1-prop} & Pointwise and singular-value extremality of the direct rotation & proved at source scope \\
\\path{DK-4.1-cor} & Unitarily invariant norm minimality of restricted displacement & proved at source scope \\
\\path{DK-4.2-prop} & Basis-angle square-sum extremality & proved at source scope \\
\\path{DK-4.3-prop} & Squared-displacement unitarily invariant norm minimality & proved at source scope \\
\\path{DK-4.4-prop} & Full-displacement claim of Proposition~4.4 & refuted as printed; repaired separately \\
\\path{DK-5.1-thm} & Banach-space Sylvester lower bound & proved at source scope \\
\\path{DK-5.2-thm} & Semibounded self-adjoint Sylvester theorem & proved at source scope \\
\\path{DK-5.1-lem} & Strong-cutoff convergence of singular values & proved at source scope \\
\\path{DK-6.1-lem} & Direct-sum norm comparison and converse & proved at source scope \\
\\path{DK-6.2-lem} & Reflection-pinch contraction & proved at source scope \\
\\path{DK-6.1-prop} & Sine proof, ambient limitation, and symmetric sine theorem & proved at source scope \\
\\path{DK-6.1-thm} & Generalized sine theorem & proved at source scope \\
\\path{DK-6.2-thm} & Pairwise-gap square-norm sine theorem & proved at source scope \\
\\path{DK-6.3-thm} & Generalized tangent theorem and supporting machinery & proved at source scope \\
\\path{DK-6.3-lem} & Finite-rank near-maximizer leakage estimate & proved at source scope \\
\\path{DK-8.1-thm} & Branch selection and spectral repulsion & proved at source scope \\
\\path{DK-8.2-thm} & Smallness selects the acute branch & proved at source scope \\
\bottomrule
\end{longtable}

The result table is keyed to source results.  One source result can require many
Lean declarations, while a single reusable Lean theorem can support more than
one source-specific wrapper.  The denominator therefore remains stable under
refactoring of implementation modules.

\section{Source-correspondence review}
\label{app:source-review}

A result is accepted only after two distinct checks.

\paragraph{Mechanical verification.}
The declaration must resolve in the ordinary project build, and every theorem
on which it depends must be accepted by Lean's kernel.  This establishes the
formal proposition exactly as encoded.

\paragraph{Source comparison.}
The encoded proposition is compared with the cited source location.  The review
checks the mathematical objects, hypotheses, conclusion, scalar field,
dimension, norm class, boundary conventions, and operator-domain assumptions.
For a nonliteral representation, the review also records the theorem or
mathematical argument that justifies the correspondence.  Equation
\cref{eq:modulus-sine} is the main example: the rectangular map in the Lean
sine-theta theorem represents the source's positive sine operator only after a
modulus/polar-decomposition argument.

This separation prevents theorem names from serving as evidence of fidelity.
For example, the proved declaration
\texttt{sinTheta\_unbounded\_intervalExterior\_characterizedWitness\_rclike}
remains a useful finite-gap theorem, but it is not the selected source endpoint
because it does not include the half-infinite gap cases.  Conversely,
Proposition~4.4 is complete in the inventory even though its printed assertion
has no proof: the completed formal treatment consists of the transcribed claim
as a proposition, a checked counterexample/refutation, a diagnosis of the
source proof, and repaired theorems.

\section{YWS source-audit scope}
\label{app:yws-coverage}

The YWS audit uses a different denominator because the journal paper does not
claim a line-by-line formalization of every mathematical statement in that
article.  The maintained census has \YWSCensusEntryCount{} rows:
\YWSCensusPaperSourceCount{} rows associated with printed source material and
\YWSCensusExplicitAdditionCount{} explicitly marked additions.  All
\YWSCensusVerifiedCount{} rows resolve to declarations in the normal build.

The source rows include the conclusions of Theorems~1--3 and the equations,
lemmas, examples, and boundary conventions needed to check those conclusions.
Several rows can share one correction.  The census has
\YWSCensusCorrectedCount{} rows whose implementation is classified as corrected,
but the manuscript reports two printed YWS defects: the missing square in
equation~(4) and the rank-boundary convention in Theorem~3.

\section{Formalization procedure and AI assistance}
\label{app:formalization-process}

The development combined ordinary Lean compilation with extensive LLM
assistance.  Work on a source theorem typically began by reading the statement
and surrounding proof, identifying standing assumptions, and locating the
relevant Lean interfaces.  The target was then decomposed into intermediate
mathematical obligations.  Existing Mathlib, Tau Ceti, local, and external Lean
results were searched before introducing new foundations.  Candidate
definitions and proofs were compiled with Lean and revised in response to
elaboration or proof failures.  A separate pass compared the accepted theorem
statement with the source, including its operator domains, scalar field,
dimension, norm class, constants, gap hypotheses, and endpoint conventions.

LLMs participated at every stage of this cycle.  They transcribed and explained
source material, proposed decompositions, searched code and literature,
implemented Lean declarations, repaired proofs from compiler feedback, and
performed adversarial source comparisons.  Human authors selected the target
scope, decided which source interpretations to accept, evaluated proposed
counterexamples and repairs, and reopened statements when later review exposed
a mismatch.  The stopping condition for a source result was therefore a
kernel-checked formal treatment together with an accepted mathematical
comparison to the source, or a checked refutation when the printed statement
was false.

The systems used during the project are summarized in
\Cref{tab:formalization-systems}.  Model assignments changed during the
project.  The table therefore records observed uses without assigning a fixed
role to any one model.

\begin{table}[H]
\centering
\small
\begin{tabularx}{\linewidth}{>{\raggedright\arraybackslash}p{0.20\linewidth}>{\raggedright\arraybackslash}p{0.28\linewidth}X}
\toprule
System & Models used & Observed uses \\
\midrule
ChatGPT chat interface & GPT-5.6 Sol & Mathematical discussion, proof planning, source comparison, review, literature search, and manuscript drafting and editing. \\
Claude Code & Claude Opus 4.8; Claude Opus 5; Claude Fable 5 & Repository search, Lean implementation, compilation, proof repair, and long-running formalization tasks. \\
Codex & GPT-5.6 Sol; GPT-5.6 Terra; GPT-5 Codex; GPT-5.6 Codex & Repository search, Lean implementation, compilation, repair, review, and proposed repository changes. \\
\bottomrule
\end{tabularx}
\caption{AI systems and model families used during the formalization and
manuscript preparation.  The prose of the paper was drafted and edited with AI
assistance; the authors reviewed and take responsibility for the submitted
text.}
\label{tab:formalization-systems}
\end{table}

Lean's kernel checks the proof terms produced by this process.  It does not
check that a theorem statement is the intended transcription of a paper.  The
statement-review procedure in Appendix~\ref{app:source-review} addresses that
second question.  Some early valid Lean statements covered narrower bounded,
finite-dimensional, or spectral-gap cases than the corresponding
Davis--Kahan theorem, so those statements were reopened after source review.

\section{Artifact organization}
\label{app:artifact}

The paper-facing Davis--Kahan declarations are under
\path{DavisKahan/Sources/DavisKahan1970}.  Reusable mathematical infrastructure
is primarily under \path{ForTauCeti}; its declarations already use final
\texttt{TauCeti.*} namespaces.  The maintained Davis--Kahan result inventory is
\path{dev/davis-kahan-1970-formalization-result-inventory.json}.  The YWS source
census is maintained separately under the YWS package and is summarized for
this manuscript by the generated census snapshot.  This organization permits
source-audit records and paper-specific counterexamples to remain local to the
paper package while reusable analysis is prepared for upstream integration.
